%! Scatterplots and Correlation — Unit 2, Lesson 2.2 — Slides (Beamer)
\documentclass[aspectratio=169, 11pt]{beamer}
\usepackage{apstats-beamer}   % provides \navyheader{} and \sectionlabel[color]{}

\begin{document}

% TITLE SLIDE (CourseName is not defined in beamer, so the name is literal)
{
\setbeamercolor{background canvas}{bg=navy}
\begin{frame}[plain]
  \vspace{2.8cm}\hspace{0.55cm}%
  \begin{minipage}{0.9\textwidth}
    {\color{goldacc}\bfseries\small LESSON 2.2}\\[0.5cm]
    {\color{white}\bfseries\LARGE Scatterplots and Correlation}\\[1.2cm]
    {\color{skymid}\small Statistics~$\cdot$~Unit 2}
  \end{minipage}
\end{frame}
}

% GRADUAL-RELEASE deck: targets → warm-up → I-DO divider → one frame per notes row (three) →
% WE-DO guided practice → spoken debrief → close & START THE HOMEWORK (the individual
% practice). There is no group activity and no independent set in the notes.
\newcommand{\redink}[1]{{\color{redacc}\bfseries #1}}

% ── LEARNING TARGETS ─────────────────────────────────────────────────────────
\begin{frame}[t, plain]
  \navyheader{Today's targets}
  \sectionlabel{Lesson 2.2}
  By the end of today, I can\ldots
  \begin{itemize}\setlength{\itemsep}{5pt}
    \item read a \textbf{scatterplot} and name its \textbf{explanatory} and
          \textbf{response} variables.
    \item describe an \textbf{association} --- direction, form, strength, unusual features ---
          in context.
    \item interpret the \textbf{correlation} $r$, and say why $r$ near 0 is not ``no
          relationship.''
    \item explain why a strong $r$ does not show \textbf{causation}.
  \end{itemize}
  \begin{block}{How today works}
    Warm-up $\to$ \textbf{notes together} $\to$ \textbf{one we work as a class} $\to$
    debrief $\to$ \textbf{you start the homework here, alone}.
  \end{block}
\end{frame}

% ── WARM-UP ──────────────────────────────────────────────────────────────────
\begin{frame}[t, plain]
  \navyheader{Warm-Up: three things you already know}
  \sectionlabel{5 minutes --- on your own}
  \begin{enumerate}\setlength{\itemsep}{7pt}
    \item 60 of 80 athletes passed a fitness test; 72 of 120 non-athletes did. Is there an
          \textbf{association}?
    \item Name the four things every description of a distribution mentions.
    \item Five students' study hours and quiz scores: as hours go up, what happens to the
          score? Which variable would you use to \emph{predict} the other?
  \end{enumerate}
  \begin{block}{Keep this on the board}
    Shape, center, spread, unusual values --- today a scatterplot gets its own four words.
  \end{block}
\end{frame}

% ── I DO — direct instruction begins ─────────────────────────────────────────
{
\setbeamercolor{background canvas}{bg=navy}
\begin{frame}[plain]
  \vspace{1.8cm}\hspace{0.8cm}%
  \begin{minipage}{0.86\textwidth}
    {\color{goldacc}\bfseries\small GUIDED NOTES \ \textbullet\ \ I DO}\\[0.7cm]
    {\color{white}\bfseries\Large Open to your Guided Notes page. Fill each blank as we
    name it --- not before.}\\[0.9cm]
    {\color{skymid}\small Three terms go in the vocabulary box today. Write each one the
    moment we use it.}
  \end{minipage}
\end{frame}
}

% ── ROW 1: TWO QUANTITATIVE VARIABLES ────────────────────────────────────────
\begin{frame}[t, plain]
  \navyheader{Two quantitative variables}
  \sectionlabel{notes, Variables $\cdot$ problems 1--3}
  \begin{columns}[T]
    \begin{column}{0.48\textwidth}
      \begin{block}{Scatterplot}
        Two quantitative variables on the same individuals --- \textbf{one dot per
        individual} (EK UNC-1.S.1, 1.S.2).
      \end{block}
      \vspace{4pt}
      \begin{block}{Explanatory $\to$ response}
        The \textbf{explanatory} variable ($x$-axis) helps explain or predict the
        \textbf{response} ($y$-axis) (EK UNC-1.S.3).
      \end{block}
    \end{column}
    \begin{column}{0.48\textwidth}
      \centering
      \begin{tikzpicture}[x=1.0cm, y=0.065cm]
        \foreach \y in {20,30,40,50} {
          \draw[linegray, very thin] (2,\y) -- (6,\y);
          \node[left, font=\tiny] at (2,\y) {\y};}
        \draw[thick] (2,15) -- (6.1,15); \draw[thick] (2,15) -- (2,52);
        \foreach \x in {2,3,4,5,6} \node[below, font=\tiny] at (\x,15) {\x};
        \foreach \w/\m in {2.5/38, 2.8/36, 3.0/35, 3.2/48, 3.3/32, 3.5/31, 3.6/29, 3.9/28,
                           4.1/26, 4.4/24, 4.8/21, 5.4/17}
          \fill[navy] (\w,\m) circle (2pt);
        \node[font=\tiny] at (4,7) {Weight (1000 lb)};
        \node[rotate=90, font=\tiny] at (1.4,33) {Highway mpg};
      \end{tikzpicture}\\[2pt]
      {\small 12 car models}
    \end{column}
  \end{columns}
\end{frame}

% ── ROW 2: DESCRIBING THE PATTERN ────────────────────────────────────────────
\begin{frame}[t, plain]
  \navyheader{Describing the pattern: DOFS}
  \sectionlabel{notes, Pattern $\cdot$ problems 4--6 $\cdot$ always in context}
  \begin{columns}[T]
    \begin{column}{0.5\textwidth}
      \begin{itemize}\setlength{\itemsep}{5pt}
        \item \textbf{D}irection --- positive or negative (EK DAT-1.A.2, 1.A.3)
        \item \textbf{O}utliers \& unusual features --- far from the pattern, or clusters
              (EK DAT-1.A.6)
        \item \textbf{F}orm --- linear or curved (EK DAT-1.A.4)
        \item \textbf{S}trength --- how closely points follow the form (EK DAT-1.A.5)
      \end{itemize}
    \end{column}
    \begin{column}{0.46\textwidth}
      \begin{block}{The cars, in one sentence}
        A strong, negative, linear association between weight and highway mpg --- heavier
        cars tend to get lower mpg. One hybrid sits far above the pattern.
      \end{block}
    \end{column}
  \end{columns}
  \vspace{10pt}
  \begin{center}
    {\large The heaviest car is far to the right --- but it \redink{fits the pattern.}}\\[4pt]
    {\small\color{slate} Unusual means far from the \emph{pattern}, not far from the middle.}
  \end{center}
\end{frame}

% ── ROW 3: READING r — THE CRUX ──────────────────────────────────────────────
\begin{frame}[t, plain]
  \navyheader{Reading the correlation $r$}
  \sectionlabel{notes, Reading $r$ $\cdot$ problems 7--10 $\cdot$ problem 9 is the crux}
  \begin{columns}[T]
    \begin{column}{0.5\textwidth}
      \begin{itemize}\setlength{\itemsep}{4pt}
        \item $r$ gives the direction and strength of the \textbf{linear} association
              (EK DAT-1.B.1).
        \item $-1 \le r \le 1$; no units; switching $x$ and $y$ leaves it alone
              (EK DAT-1.C.1).
        \item Cars: $r = -0.86$ --- negative and fairly strong.
        \item Correlation is \textbf{not} causation (EK DAT-1.C.2).
      \end{itemize}
    \end{column}
    \begin{column}{0.46\textwidth}
      \centering
      \begin{tikzpicture}[x=0.068cm, y=0.16cm]
        \foreach \y in {20,25,30,35} {
          \draw[linegray, very thin] (15,\y) -- (80,\y);
          \node[left, font=\tiny] at (15,\y) {\y};}
        \draw[thick] (15,20) -- (81,20); \draw[thick] (15,20) -- (15,37);
        \foreach \x in {20,40,60,80} \node[below, font=\tiny] at (\x,20) {\x};
        \foreach \s/\m in {20/23, 25/27, 30/30, 35/32, 40/34, 45/35, 50/35, 55/34, 60/32,
                           65/30, 70/27, 75/24}
          \fill[navy] (\s,\m) circle (2pt);
        \node[font=\tiny] at (47.5,16.5) {Speed (mph)};
        \node[rotate=90, font=\tiny] at (7,28) {mpg};
      \end{tikzpicture}\\[2pt]
      {\small One car, 12 speeds: $r = 0.03$}
    \end{column}
  \end{columns}
  \vspace{8pt}
  \begin{center}
    {\large $r = 0.03$ --- and yet speed \redink{clearly matters.}}\\[4pt]
    {\small\color{slate} $r$ near 0 means no \emph{linear} association. Look at the plot first.}
  \end{center}
\end{frame}

% ── WE DO ────────────────────────────────────────────────────────────────────
\begin{frame}[t, plain]
  \navyheader{Guided Practice: The City Pool}
  \sectionlabel[goldacc]{We do $\cdot$ you hold the pen $\cdot$ problems 11--13}
  High temperature and visitors at a city pool on 14 summer days; $r = 0.90$.
  \begin{itemize}\setlength{\itemsep}{6pt}
    \item Which variable is explanatory, and which is the response?
    \item Describe the association in context --- direction, form, strength, unusual points,
          and $r$.
    \item Ice-cream sales and visitors have $r = 0.88$. Does ice cream bring swimmers?
  \end{itemize}
  \begin{block}{The question behind the question}
    When two things rise together, ask: \redink{what else rises with both?}
  \end{block}
\end{frame}

% ── DEBRIEF ──────────────────────────────────────────────────────────────────
\begin{frame}[t, plain]
  \navyheader{Debrief: what $r$ can and cannot say}
  \sectionlabel{8 minutes $\cdot$ pens down, papers up --- we say these aloud}
  \vspace{6pt}
  \begin{columns}[T]
    \begin{column}{0.46\textwidth}
      \begin{block}{$r$ near 0}
        No \emph{linear} association. The speed plot is a strong curve with $r = 0.03$.
      \end{block}
    \end{column}
    \begin{column}{0.46\textwidth}
      \begin{block}{$r$ near $\pm 1$}
        A strong linear pattern --- still not proof that $x$ causes $y$.
      \end{block}
    \end{column}
  \end{columns}
  \vspace{14pt}
  \begin{block}{Say it out loud}
    Name something $r$ tells you, and something only the \redink{picture} can tell you.
  \end{block}
\end{frame}

% ── CLOSE & START THE HOMEWORK ───────────────────────────────────────────────
\begin{frame}[t, plain]
  \navyheader{Homework --- start it now}
  \sectionlabel[goldacc]{Scored $\cdot$ due the first class after two study halls}
  \begin{itemize}\setlength{\itemsep}{5pt}
    \item \textbf{Part A:} practice hours and free throws --- variables, then a description
          in context.
    \item \textbf{Part B:} the strongest $r$, a U-shaped plot with $r = 0.05$, and libraries
          and crime.
    \item \textbf{AP Practice --- not scored:} screens and sleep. Work it for the reps.
  \end{itemize}
  \begin{block}{Silent and alone --- I am circulating}
    \emph{Part B, item 2} is the one I am reading over your shoulder. If you wrote ``no
    relationship,'' flag me before you leave.
  \end{block}
\end{frame}

% ── CLOSE ────────────────────────────────────────────────────────────────────
{
\setbeamercolor{background canvas}{bg=navy}
\begin{frame}[plain]
  \vspace{1.5cm}\hspace{0.8cm}%
  \begin{minipage}{0.86\textwidth}
    {\color{goldacc}\bfseries\small BEFORE YOU GO}\\[0.7cm]
    {\color{white}\bfseries\Large $r$ measures a straight line. Look at the picture before
    you believe the number.}\\[0.9cm]
    {\color{skymid}\small Next: Lesson 2.3 --- Regression Lines and Residuals. When the
    pattern is linear, we draw the line and use it to predict.}
  \end{minipage}
\end{frame}
}

\end{document}
